\chapter{系统测试}
\label{cha:system_test}

\section{测试目标与范围}

第五章不再沿着实现模块逐项罗列接口现象，而是直接回到第二章提出的功能需求与非功能需求，考察系统在真实运行环境中是否已经形成可持续使用的业务闭环。这里所关心的并不是某个局部函数是否返回预期值，而是身份建立、文档接入、检索过滤、文件服务和会话交互这些能力能否在同一套系统中连续成立。

因此，本章的判断边界也是明确的。测试结论仅用于说明系统在数据库、缓存、对象存储、消息中间件、检索引擎与外部模型服务共同参与的条件下，已经达到怎样的运行状态；它不用于证明算法效果优于其他系统，也不把个别问答样例直接外推为开放场景下的通用最优表现。

\section{测试环境与需求映射}

本轮测试采用真实依赖环境，并将第二章需求分析中的关键要求重新归并为可观测的验证项。环境与数据规模见表 \ref{tab:system-test-env}，需求映射结果见表 \ref{tab:function-requirement-result} 和表 \ref{tab:nonfunction-requirement-result}。

\begin{table}[H]
    \centering
    \caption{系统测试环境与数据规模}
    \label{tab:system-test-env}
    {\small
    \begin{tabularx}{\textwidth}{L{3.2cm}Y}
        \toprule
        项目 & 说明 \\
        \midrule
        操作系统 & Linux \\
        Java 版本 & 17.0.18 \\
        运行依赖 & MySQL、Redis、MinIO、Kafka、Elasticsearch \\
        外部模型服务 & \texttt{Qwen text-embedding-v4}（2048 维）与 \texttt{deepseek-chat} \\
        检索索引 & \texttt{knowledge\_base} \\
        消息主题 & \texttt{file-processing-system} 及其死信主题 \\
        文档总量 & 500 份 \\
        全链路入库样本 & 100 份 \\
        仅上传样本 & 400 份 \\
        大文档样本 & 100 份 \\
        检索查询数 & 40 条 \\
        权限矩阵检查数 & 25 组 \\
        文件服务检查数 & 37 项 \\
        流式交互轮次 & 5 轮 \\
        固定分片大小 & 5 MB \\
        \bottomrule
    \end{tabularx}
    }
\end{table}

\begin{table}[H]
    \centering
    \caption{功能需求验证结果}
    \label{tab:function-requirement-result}
    {\small
    \begin{tabularx}{\textwidth}{L{1.2cm}L{3.9cm}Y L{1.8cm}}
        \toprule
        编号 & 需求内容 & 核心结果 & 结论 \\
        \midrule
        F1 & 用户注册、登录后形成可持续身份上下文 & 注册与登录检查通过率 100\% & 满足 \\
        F2 & 私人组织标签自动创建，支持组织分配与主组织切换 & 8 项身份与组织检查全部通过 & 满足 \\
        F3 & 支持文档上传、续传、合并、解析与知识入库 & 上传、续传、合并一致率均为 100\%，解析与索引通过率为 99\% & 基本满足 \\
        F4 & 支持自然语言混合检索并返回可用候选结果 & Recall@5 为 1.0000，MRR@5 为 0.9813 & 满足 \\
        F5 & 权限过滤在检索阶段完成 & 25 组权限矩阵全部通过，越权返回率为 0 & 满足 \\
        F6 & 支持文件列表、预览、下载和删除 & 37 项文件服务检查全部通过 & 满足 \\
        F7 & 支持流式回答、多轮上下文与历史记录读取 & 流式输出与历史读取成立，多轮历史追问未通过 & 未完全满足 \\
        \bottomrule
    \end{tabularx}
    }
\end{table}

\begin{table}[H]
    \centering
    \caption{非功能需求验证结果}
    \label{tab:nonfunction-requirement-result}
    {\small
    \begin{tabularx}{\textwidth}{L{1.2cm}L{3.9cm}Y L{1.8cm}}
        \toprule
        编号 & 需求内容 & 核心结果 & 结论 \\
        \midrule
        NF1 & 离线入库与在线交互具备可量化时延指标 & 上传、可检索时延、检索时延、TTFT 与完整响应时间均可统计 & 满足 \\
        NF2 & 检索与文件服务全过程不得越权 & 检索越权返回率为 0，文件服务检查通过率为 100\% & 满足 \\
        NF3 & 异常下保持可恢复与可持续推进 & 续传成功率 100\%，完成通知到达率 100\%，历史读取可用 & 满足 \\
        NF4 & 系统结构保持分层清晰、配置解耦、组件可替换 & 6 项结构符合性检查全部通过 & 满足 \\
        \bottomrule
    \end{tabularx}
    }
\end{table}

\section{功能需求验证}

从结果看，本轮功能测试更接近对业务闭环的核查，而不是对单个接口的点验。身份建立决定后续上传和检索是否具有持续的权限上下文；文档接入决定知识是否能够进入可检索状态；检索、文件服务与流式问答则对应用户真正消费知识的过程。只有把这些环节放在一起看，第二章提出的功能需求才算得到有效验证。

\subsection{身份建立与组织管理}

身份与组织管理部分的 8 项检查全部通过。用户完成注册后即可获得私人组织标签，登录后能够读取当前身份信息；管理员可为用户创建并分配组织标签，用户切换主组织后，新的组织上下文能够继续被上传场景读取。换言之，系统中的“用户身份”并不是一次性认证结果，而是可以持续参与后续业务处理的权限上下文。这一结果意味着第二章关于用户注册、组织归属和主组织切换的要求已经落到实际运行层面。

这一部分虽然不直接生成知识片段，也不直接返回问答结果，却决定了后续所有业务动作的边界。若身份与组织关系不能稳定传递，上传归属、检索过滤和文件服务都会失去依据。从本轮结果看，这一基础前提已经成立。

\subsection{文档接入与知识入库}

文档接入与知识入库的关键结果见表 \ref{tab:document-processing-result}。

\begin{table}[H]
    \centering
    \caption{文档接入与知识入库结果}
    \label{tab:document-processing-result}
    {\small
    \begin{tabularx}{\textwidth}{L{4.2cm}L{2.5cm}Y}
        \toprule
        指标 & 结果 & 说明 \\
        \midrule
        样本规模 & 500 份 & 其中 100 份进入完整入库链路，400 份用于上传与续传验证 \\
        文件类型校验通过率 & 100\% & 所有测试文档均通过格式校验 \\
        上传成功率 & 100\% & 500 份文档均完成上传 \\
        断点续传成功率 & 100\% & 119 个续传样本全部恢复成功 \\
        合并一致率 & 100\% & 所有合并后对象的 MD5 与原文件一致 \\
        解析成功率 & 99\% & 100 份全链路样本中有 99 份完成解析 \\
        索引写入成功率 & 99\% & 与解析结果保持一致 \\
        存储一致性通过率 & 99\% & 数据库、对象存储与检索索引三端一致 \\
        平均上传延迟 & 71.182 ms & 单文档上传阶段平均耗时 \\
        P95 上传延迟 & 148.0 ms & 上传阶段长尾耗时 \\
        平均合并延迟 & 65.672 ms & 服务端合并平均耗时 \\
        P95 合并延迟 & 153.0 ms & 合并阶段长尾耗时 \\
        平均可检索就绪延迟 & 94362.46 ms & 从合并完成到知识可检索的平均耗时 \\
        P95 可检索就绪延迟 & 142574.0 ms & 全链路入库长尾耗时 \\
        \bottomrule
    \end{tabularx}
    }
\end{table}

可以看到，文档接入链路在传输层表现得较为稳定。500 份样本全部完成上传，119 个断点续传样本全部恢复成功，所有合并后对象的 MD5 与原文件一致。这说明第二章提出的上传、续传和服务端合并需求已经得到直接验证，传输阶段没有出现状态分裂或对象损坏。

真正需要谨慎表述的是入库结果。在 100 份进入完整链路的文档中，有 99 份完成了解析、向量化和索引写入，仍有 1 份文档未形成向量记录和检索索引。也就是说，当前系统已经能够完成大多数文档的稳定入库，但这一能力尚不能表述为“全部样本无异常完成”。就需求结论而言，文档接入与知识入库需求已经基本满足，不过文档解析稳定性仍有继续补强的必要。

\subsection{检索、权限与文件服务}

知识消费相关功能并不只体现在“能否检索到文档”这一点上。对企业文档系统而言，候选结果是否相关、候选集合是否受权限约束，以及用户是否能够继续完成预览、下载和删除，三者应当被放在同一组需求中观察。相关结果见表 \ref{tab:retrieval-permission-file-result}。

\begin{table}[H]
    \centering
    \caption{检索、权限与文件服务结果}
    \label{tab:retrieval-permission-file-result}
    {\small
    \begin{tabularx}{\textwidth}{L{3.2cm}L{3.0cm}Y}
        \toprule
        验证对象 & 结果 & 说明 \\
        \midrule
        混合检索相关性 & Recall@1=0.9750，Recall@5=1.0000 & 40 条查询全部在 Top5 命中目标文档，其中 39 条首位命中 \\
        混合检索排序质量 & MRR@5=0.9813，nDCG@5=0.9858 & 排序总体稳定，但仍存在个别非首位命中样例 \\
        检索响应时间 & 平均 357.7 ms，P95 836.0 ms & 检索阶段时延保持在亚秒级到秒级区间 \\
        权限矩阵检查 & 25/25 通过 & 5 类用户与 5 类文档的访问关系全部符合预期 \\
        检索越权返回率 & 0 & 125 条检索结果中未出现未授权文档 \\
        文件服务检查 & 37/37 通过 & 列表、预览、下载和删除均遵守权限边界 \\
        删除后状态清理 & 通过 & 数据库、对象存储、向量记录与检索索引同步清除 \\
        \bottomrule
    \end{tabularx}
    }
\end{table}

从检索结果看，系统已经能够为问答链路提供稳定的候选证据。40 条查询全部在前 5 条结果中命中目标文档，其中只有 1 条没有排在首位，而是落在第 4 位。这说明当前混合检索已经不再停留在“偶尔能找回目标文档”的阶段，而是具备了较高的相关性稳定度。不过，唯一的非首位命中样例也提示出一个事实：当查询对象处于相邻年份、标题高度相似的文档群中时，排序仍会受到语义近邻干扰。就需求而言，F4 已经满足；就质量边界而言，排序细节仍未完全收敛。

权限与文件服务部分的结果更为整齐。25 组访问矩阵检查全部通过，检索阶段越权返回率为 0，说明权限约束已经真正进入候选结果的构造过程，而不是停留在接口入口的形式性判断。与此相对应，文件列表、预览、下载和删除共 37 项检查全部通过，表明权限边界在文件消费环节同样成立。换句话说，系统已经把“谁可以看到什么、谁可以取走什么、谁可以删除什么”落实到了完整的使用链路中。

\subsection{流式问答与历史记录}

流式交互结果见表 \ref{tab:streaming-interaction-result}。

\begin{table}[H]
    \centering
    \caption{流式问答与历史记录结果}
    \label{tab:streaming-interaction-result}
    {\small
    \begin{tabularx}{\textwidth}{L{4.2cm}L{2.5cm}Y}
        \toprule
        指标 & 结果 & 说明 \\
        \midrule
        交互轮次 & 5 轮 & 包含文档问答、历史追问和后续问答 \\
        连接建立情况 & 成功 & WebSocket 连接能够正常建立 \\
        完成通知到达率 & 100\% & 每轮均收到完成消息 \\
        流式轮次通过率 & 100\% & 每轮均产生连续分块输出 \\
        平均 TTFT & 1302.6 ms & 首个分块平均返回时间 \\
        P95 TTFT & 1562.0 ms & 首字等待时间长尾波动 \\
        平均完整响应时间 & 13772.8 ms & 从提问到完成通知的平均耗时 \\
        P95 完整响应时间 & 19615.0 ms & 完整响应长尾耗时 \\
        平均分块数 & 223.8 & 外部模型条件下保持连续流式输出 \\
        平均块间隔 & 30.5963 ms & 相邻分块之间的平均间隔 \\
        历史记录读取 & 通过 & 历史接口返回 10 条消息，耗时 17 ms \\
        多轮历史追问 & 未通过 & 两次追问均出现语义偏移 \\
        \bottomrule
    \end{tabularx}
    }
\end{table}

从交互形态上看，系统已经具备稳定的流式回答能力。五轮交互均成功建立连接，均能够持续输出分块，并在结束时返回完成通知；三轮基于具体文档的问答也都给出了带依据的回答。这说明流式响应、引用组织和会话消息发送机制已经进入可用状态，F7 中关于“能够进行流式回答”的部分是成立的。

但同一项需求并未整体完成。针对“上一轮提到的公司全称是什么”“上一轮对应的文档名称是什么”这两次追问，系统返回的内容与首轮问题并不一致，历史语义没有稳定延续到后续生成过程。换言之，历史记录已经能够被保存和读取，却还没有在多轮追问中形成可靠的上下文约束。因此，本章对 F7 的判断只能是“未完全满足”，而不能笼统写成“多轮会话能力已经稳定成立”。

\section{非功能需求验证}

非功能需求并不表现为某一条接口是否存在，而是体现在系统运行后的边界上。对企业文档场景而言，这些边界主要落在时延、安全、可靠性和可维护性四个方面。它们不决定系统有没有功能，却决定这些功能能否被持续使用。

\begin{table}[H]
    \centering
    \caption{主要性能指标结果}
    \label{tab:performance-result}
    {\small
    \begin{tabularx}{\textwidth}{L{3.5cm}L{3.2cm}Y}
        \toprule
        性能对象 & 结果 & 说明 \\
        \midrule
        上传阶段时延 & 平均 71.182 ms，P95 148.0 ms & 单文档上传处理耗时 \\
        合并阶段时延 & 平均 65.672 ms，P95 153.0 ms & 服务端对象合并耗时 \\
        可检索就绪时延 & 平均 94362.46 ms，P95 142574.0 ms & 从上传完成到知识可检索的全链路耗时 \\
        检索响应时延 & 平均 357.7 ms，P95 836.0 ms & 在线检索阶段响应时间 \\
        首字等待时间 & 平均 1302.6 ms，P95 1562.0 ms & 流式问答首个分块返回时间 \\
        完整响应时间 & 平均 13772.8 ms，P95 19615.0 ms & 从提问到回答完成的总耗时 \\
        历史读取时延 & 17 ms & 历史消息读取接口耗时 \\
        \bottomrule
    \end{tabularx}
    }
\end{table}

从性能维度看，系统已经具备可量化的在线与离线时延画像。检索阶段平均耗时为 357.7 ms，说明候选构造本身并未成为主要瓶颈；流式问答能够在约 1.3 s 内返回首个分块，也说明在线交互已经具备基本可接受的反馈节奏。相比之下，离线入库阶段的主要耗时集中在解析、向量化和索引可见性等待，知识从上传完成到进入可检索状态平均需要 94.36 s。这一结果并不影响“系统能够完成异步入库”这一需求结论，但它同时表明，当前版本更适合后台入库场景，而不适合把“上传后立即问答”作为默认体验。

\begin{table}[H]
    \centering
    \caption{安全性、可靠性与可维护性结果}
    \label{tab:security-reliability-maintainability-result}
    {\small
    \begin{tabularx}{\textwidth}{L{3.5cm}L{3.2cm}Y}
        \toprule
        需求属性 & 结果 & 说明 \\
        \midrule
        检索安全性 & 越权返回率 0 & 权限过滤已实际作用于候选结果生成 \\
        文件服务安全性 & 37 项检查全部通过 & 预览、下载和删除均未出现越权访问 \\
        传输恢复能力 & 119 个续传样本全部成功 & 网络中断后可恢复未完成上传 \\
        状态一致性 & 合并一致率 100\%，存储一致性 99\% & 传输稳定，入库侧仍有个别异常样本 \\
        完成通知能力 & 到达率 100\% & 流式交互结束状态能够被可靠感知 \\
        历史读取能力 & 通过 & 历史记录可以稳定返回 \\
        结构符合性 & 6 项检查全部通过 & 分层结构、配置解耦、降级路径和组件替换空间均已具备 \\
        \bottomrule
    \end{tabularx}
    }
\end{table}

安全性方面，本轮测试给出的结果比较明确：无论是检索结果还是文件服务，都没有出现越权返回。对企业文档系统来说，这一点比单纯的召回率更重要，因为一旦候选结果跨越组织边界，后续生成能力越强，风险也会随之扩大。本轮结果至少说明，当前版本已经把访问边界落实到了知识消费过程，而不只是停留在登录态检查。

可靠性与可维护性则表现出另一种特征。系统已经具备断点恢复、完成通知和历史读取等基础能力，分层结构、配置解耦、异常降级与死信处理等设计也没有出现明显缺口；但与之相伴的现实情况是，入库侧仍存在个别样本未完成解析与索引写入，多轮历史追问也尚未稳定成立。也就是说，当前版本的系统结构已经为后续扩展和修正留出了空间，但这并不意味着所有运行细节都已完全收敛。

\section{本章小结}

将第二章需求与本章结果对照起来看，F1、F2、F4、F5、F6 以及 NF1、NF2、NF3、NF4 均获得了直接验证，F3 达到基本满足，F7 则未完全满足。由此可以判断，系统已经形成了从身份建立、文档接入、权限受控检索到文件消费的主要业务闭环，并具备了继续进入实际使用场景的基础条件。

同时，本章也保留了两个不能被弱化处理的边界。其一，知识入库仍有 1 份全链路样本未完成解析与索引写入，说明文档解析与入库稳定性仍需继续补强；其二，流式问答虽然能够稳定建立连接并持续返回分块，但多轮历史追问出现了明显语义偏移，会话连续性尚未达到稳定状态。系统测试因此得出的不是“系统已经全部完善”的结论，而是“主要需求已经成立，剩余问题已经明确”。
